\chapter{Lean 4 for Physicists: A Crash Course}\label{ch:leanprimer}

Every earlier chapter ends with a section ``What the machine has checked'', and every
theorem there carries the badge \TA. This chapter explains what that badge costs and what it
buys. The question it answers is practical: \emph{how does a physicist read a Lean 4
statement, decide whether it says what the prose claims, and check it on her own machine?}
The answer does not require learning to write proofs from scratch. It requires three things:
knowing what a type is and why a theorem is one; knowing how Mathlib writes a matrix, a
transpose and an inverse; and knowing what the dozen tactics that appear in this book's
proofs actually do. We build these up on the four files that the rest of the book quotes
most often---the lattice definitions of \cref{ch:lattices}, the duality group of
\cref{ch:odd}, the dual-scale bound of \cref{ch:dualscale} and the circle mass spectrum of
\cref{ch:tduality}---and then read one proof, \lean{thetaShift\_isODD}, line by line. A
student should care because a proof assistant is only as useful as the reader who can tell
a strong theorem from a weak one: the kernel checks that a proof proves its statement, and
nothing else. Reading the statement is the reader's job, and this chapter teaches it.

%=====================================================================================
\section{What a proof assistant is, and what it is not}\label{sec:leanprimer:kernel}

Lean 4\index{Lean 4} is a programming language whose type system is strong enough to
express mathematical statements, together with a small program, the
\emph{kernel}\index{kernel (Lean)}, that checks whether a given term has a given type. A
theorem is a type; a proof is a term of that type; ``the theorem is proved'' means the
kernel has accepted the term. Everything else---the tactic language, the editor, the
automation, the large-language-model agents described in \cref{ch:pipeline}---is
machinery for \emph{producing} the term. None of it has to be trusted. If a tactic has a
bug and produces a wrong term, the kernel rejects it.

Three things are trusted, and the student should know them by name.
\begin{enumerate}[leftmargin=*,itemsep=1pt]
\item \textbf{The kernel itself}, a small program by the standards of a compiler, small
  enough that independent re-implementations exist and can be run on the same proof terms.
\item \textbf{The axioms.}\index{axioms (Lean)} Lean's logic is a dependent type theory;
  its core library declares exactly three axioms: \lean{propext} (two logically equivalent propositions are
  equal), \lean{Quot.sound} (quotients behave as quotients) and \lean{Classical.choice}
  (a choice function exists, which gives the law of excluded middle). Every theorem in this
  book depends on at most these three. The command \lean{\#print axioms} lists what a
  theorem depends on (\cref{sec:leanprimer:axioms}); anything beyond the three is a
  warning sign.
\item \textbf{The statement.} The kernel checks that the proof proves the statement as
  written. It does not check that the statement is the one you meant, that the definition
  \lean{leftMomentum} is what a physicist calls the left-moving momentum, or that a matrix
  named \lean{cartanE8} is the Cartan matrix of $E_8$. This is why the bible of this book
  says that the identification of a Lean object with a physical system is never Tier A.
\end{enumerate}

\begin{pitfallbox}[The kernel checks the proof, not the statement]
Two incidents from this project's own log illustrate the third point.
A revision brief said that the $\Theta$-shift preserves $\eta$ ``iff $\Theta$ is
antisymmetric''; the Lean theorem proves only the ``if'' direction
\source{LL.md §S2.13, ll.~109--112}. And since the kernel cannot notice that a hypothesis
was quietly weakened between review and proof, the project hashes every theorem statement
into a lock file\index{statement lock} and re-checks the hashes before accepting a proof
\source{LL.md §S2.10, ll.~87--92}. The lesson for a reader is simple: before quoting a
theorem, read its statement in the \lean{.lean} file, not the prose that summarizes it.
\end{pitfallbox}

\begin{physicsbox}[What ``formalizing physics'' can mean]
No proof assistant checks that nature obeys an equation. What can be checked is a chain of
implications: \emph{if} the mass formula is $M^2 = p_L^2 + p_R^2 + \dots$ and \emph{if}
T-duality acts as $(n,w,R)\mapsto(w,n,\ap/R)$, \emph{then} $M^2$ is invariant. The two
``if''s are definitions in a file; the ``then'' is a theorem. Older libraries in this
repository go further and replace real numbers by integers or rationals: those theorems are
\emph{arithmetic shadows}\index{arithmetic shadow} of the physical statement, true as
arithmetic, silent about the continuum. The four files of this chapter are honest about
which kind they are: \lean{TDualityMassSpectrum.lean} works over $\Q$ with $\ap=1$;
\lean{TraceBound.lean} works over $\R$ with genuine real matrices.
\end{physicsbox}

%=====================================================================================
\section{Terms, types and propositions}\label{sec:leanprimer:types}

Every expression in Lean has a type, written with a colon: \lean{(3 : ℕ)}, \lean{(R : ℝ)},
\lean{(G : Matrix (Fin d) (Fin d) ℝ)}. Types themselves have types: \lean{ℝ : Type}. The
key idea, due to Curry and Howard and usually called \emph{propositions as
types}\index{propositions as types}, is that a mathematical statement is also a type, whose
inhabitants are its proofs. Statements live in a special universe called
\lean{Prop}\index{Prop@\lean{Prop}}. So
\begin{leancode}
2 ≤ R + R⁻¹          -- a Prop: the type of proofs that R + 1/R ≥ 2
Even (G i i)         -- a Prop: the type of proofs that G i i is even
Gᵀ = G               -- a Prop
\end{leancode}
are types, and a theorem is a named term of such a type. \Cref{tab:leanprimer:dict}
translates the vocabulary.

\begin{table}[ht]\centering\footnotesize
\begin{tabular}{@{}p{2.3cm}p{4.6cm}p{7.6cm}@{}}\toprule
Mathematics & Lean & Example from the repository \\ \midrule
a set or a space & a \lean{Type} & \lean{Matrix (Fin d) (Fin d) ℤ} \\
an element & a term \lean{x : T} & \lean{eta d : Matrix (Charge d) (Charge d) ℤ} \\
a statement & a term of \lean{Prop} & \lean{IsODD g} \\
a proof & a term of that \lean{Prop} & \lean{eta\_isODD : IsODD (eta d)} \\
a definition & \lean{def} & \lean{def leftMomentum (n w R : ℚ) : ℚ := n / R + w * R} \\
a theorem & \lean{theorem} & \lean{theorem eta\_mul\_self : eta d * eta d = 1} \\
``for all $x$'' & \lean{∀ x, ...} or an argument \lean{(x : T)} & \lean{∀ i, Even (G i i)} \\
``there is $x$'' & \lean{∃ x, ...} & \lean{∃ i : Fin 26, M24RepDim i = ...} \\
``and'' / ``or'' & \lean{∧} / \lean{∨} & \lean{G.det = 1 ∨ G.det = -1} \\
\bottomrule\end{tabular}
\caption{The dictionary between ordinary mathematics and Lean. All examples are taken
from the files quoted in this chapter.}\label{tab:leanprimer:dict}
\end{table}

\subsection{\lean{def} versus \lean{theorem}}
A \lean{def} introduces an object that later code can \emph{compute with} and
\emph{unfold}. A \lean{theorem} introduces a proof, and Lean deliberately forgets the
proof's internal structure once it is checked (``proof irrelevance''): two proofs of the
same proposition are equal, and a theorem can never be unfolded, only \emph{applied}. Here
is the circle mass spectrum as it is written in
\lean{StringTheoryFormalization/UseCases/TDualityMassSpectrum.lean}:
\begin{leancode}
def leftMomentum (n w R : ℚ) : ℚ := n / R + w * R
def rightMomentum (n w R : ℚ) : ℚ := n / R - w * R
def momentumMassSq (n w R : ℚ) : ℚ :=
  (leftMomentum n w R) ^ 2 + (rightMomentum n w R) ^ 2

theorem tduality_fixes_left_momentum (n w R : ℚ) (hR : R ≠ 0) :
    leftMomentum w n (1 / R) = leftMomentum n w R := by ...
\end{leancode}
Read the theorem as a sentence: ``for all rationals $n,w,R$, and given a proof \lean{hR}
that $R\neq0$, the left momentum of the dual data $(w,n,1/R)$ equals the left momentum of
$(n,w,R)$.'' The arguments before the colon are universally quantified; \lean{hR} is a
hypothesis, itself a term whose type is the proposition \lean{R ≠ 0}. The conventions are
those of \cref{ch:tduality} with $\ap=1$: $p_L = n/R + wR$, $p_R = n/R - wR$, and the file's
header says so \source{TDualityMassSpectrum.lean header, ll.~17--22}. Note what is
\emph{not} in the statement: no oscillators, no $\ap$, no real numbers. The theorem is
about rational functions of $R$; it is an arithmetic shadow of Tong's statement, and a
faithful one, because the algebra is the same over $\Q$ and $\R$.

\subsection{Implicit arguments and \lean{variable}}
Curly braces mark an argument Lean should infer: in \lean{IsEvenDiag \{n : ℕ\} (G : Gram n)}
the rank \lean{n} is read off from the type of \lean{G}, so the user writes
\lean{IsEvenDiag G}, not \lean{IsEvenDiag 8 G}. A line \lean{variable \{d : ℕ\}} at the
top of a file makes \lean{d} an implicit argument of every declaration below that mentions
it; that is why the theorems of \lean{ODD.lean} are stated for arbitrary $d$ without
repeating \lean{(d : ℕ)} each time.

\subsection{Reading a hypothesis as a physical assumption}
Hypotheses carry physics. In \lean{dualScale\_ge} the hypothesis \lean{hG : G.PosDef}
says that the torus metric is a genuine Riemannian metric; in \lean{dualScale\_inv} the
weaker \lean{hG : IsUnit G.det} says only that $G$ is invertible, because T-duality
$G\mapsto G^{-1}$ makes sense for any invertible $G$, positive or not. A theorem with a
weaker hypothesis is a stronger theorem. A theorem with no hypothesis at all, such as
\lean{eta\_mul\_self}, is a pure identity. Get into the habit of reading the hypothesis
list first: it is the fastest way to see what a formal theorem does and does not cover.

%=====================================================================================
\section{Definitions and structures}\label{sec:leanprimer:structures}

\subsection{Predicates}
The lattice file \lean{DualScaleStream2/Lattice/Basic.lean} contains the two definitions on
which every ``this lattice is even and unimodular'' statement of \cref{ch:lattices,ch:k3lattice}
rests:
\begin{leancode}
abbrev Gram (n : ℕ) := Matrix (Fin n) (Fin n) ℤ
def IsEvenDiag {n : ℕ} (G : Gram n) : Prop := ∀ i, Even (G i i)
def IsUnimodular {n : ℕ} (G : Gram n) : Prop := G.det = 1 ∨ G.det = -1
\end{leancode}
\lean{abbrev} is a \lean{def} that Lean unfolds eagerly: a \lean{Gram n} \emph{is} an
$n\times n$ integer matrix, and every matrix lemma applies to it unchanged. The other two
are \emph{predicates}\index{predicate}: functions from matrices to \lean{Prop}. Note the
honesty of the names. \lean{IsEvenDiag} asks only that the diagonal be even; that this
implies evenness of the whole quadratic form $x^{\mathsf T}Gx$ is a separate theorem,
\lean{even\_quadratic\_form\_of\_even\_diag}, which needs the extra hypothesis
\lean{Gᵀ = G}. A lattice is identified with its Gram matrix in a fixed basis; no abstract
$\Z$-module is introduced, and the file header says so
\source{Lattice/Basic.lean header, ll.~29--38}.

\subsection{Structures}
A \lean{structure}\index{structure (Lean)} bundles named fields, like a C \texttt{struct} or a Python dataclass.
From the file \lean{K3T2Signature.lean} in the directory
\lean{DualScaleStream2/Lattice/}:
\begin{leancode}
structure Signature where
  pos : ℕ
  neg : ℕ
  deriving DecidableEq, Repr

instance : Add Signature := ⟨fun a b => ⟨a.pos + b.pos, a.neg + b.neg⟩⟩
@[simp] theorem add_pos (a b : Signature) : (a + b).pos = a.pos + b.pos := rfl
\end{leancode}
A term of type \lean{Signature} is built with angle brackets, \lean{⟨3, 19⟩}, and its
fields are read with dot notation, \lean{s.pos}. The \lean{instance} line teaches Lean what
\lean{+} means on signatures (the direct sum adds the positive and negative dimensions
separately), and the theorem \lean{add\_pos} is proved by \lean{rfl}, ``reflexivity'':
both sides unfold to the same term. The attribute \lean{@[simp]} registers it as a rewrite
rule for the \lean{simp} tactic. Notice, again, what the structure does not contain: a
\lean{Signature} is a pair of naturals with no attached quadratic form. The statement
``the K3 lattice has signature $(3,19)$'' is, in this file, a \lean{def} assigning the pair
$\langle 3,19\rangle$ to a name; the theorem \lean{sigK3\_matches\_hodge} then checks that
this pair agrees with the Hodge-theoretic numbers tabulated independently in Stream~1. That
is a consistency check between two tables, not a derivation of either.

\begin{pitfallbox}[Multi-name structure fields: parenthesize them]
In Lean 4 the declaration
\begin{leancode}
structure Foo where
  a b c d : Nat        -- WRONG: one field `a`, a function of b c d
\end{leancode}
does not declare four natural-number fields. It declares a \emph{single} field \lean{a}
whose type is a curried function taking auto-bound implicit arguments \lean{b c d} and
returning \lean{Nat}. Any later use of \lean{a}, \lean{b}, \lean{c}, \lean{d} as scalars
fails with dependent-function-type errors that do not mention the cause. This reproduces
in vanilla Lean 4 with no Mathlib. The fix is to parenthesize the group,
\lean{(a b c d : Nat)}, or to write one field per line as \lean{Signature} does. The bug
was found twice, independently, in this repository's Stream~1 code
\source{LL.md §0.4, ll.~188--199}.
\end{pitfallbox}

%=====================================================================================
\section{Matrices in Mathlib}\label{sec:leanprimer:matrices}

All the block-matrix algebra of \cref{ch:odd,ch:genmetric,ch:dualscale} is written against
Mathlib's \lean{Matrix}. A few facts make the notation readable.

\paragraph{A matrix is a function.}\index{Matrix@\lean{Matrix} (Mathlib)}
\index{fromBlocks@\lean{fromBlocks}} \lean{Matrix m n α} is, by definition, the function
type \lean{m → n → α}: a matrix is a function of two indices. Hence \lean{G i j} is the
entry, \lean{fun i j => if i < j then G i j else 0} is the strict upper triangle, and two
matrices are equal when they agree at every pair of indices (the \lean{ext} tactic,
\cref{sec:leanprimer:tactics}). The index types need not be \lean{Fin n}: any finite type
will do, and \lean{ODD.lean} indexes charge vectors by the disjoint union
\begin{leancode}
abbrev Charge (d : ℕ) := Fin d ⊕ Fin d
\end{leancode}
so that a $2d\times2d$ matrix is naturally a $2\times2$ array of $d\times d$ blocks.
\index{Fin d@\lean{Fin d}}

\paragraph{Blocks.} \lean{fromBlocks A B C D} is the block matrix
$\begin{psmallmatrix}A&B\\C&D\end{psmallmatrix}$ on \lean{Fin d ⊕ Fin d}. The two Mathlib
lemmas that drive every proof in \lean{ODD.lean} are
\begin{align}
(\mathrm{fromBlocks}\,A\,B\,C\,D)^{\mathsf T} &= \mathrm{fromBlocks}\,A^{\mathsf T}\,C^{\mathsf T}\,B^{\mathsf T}\,D^{\mathsf T},
\label{eq:leanprimer:fbT}\\
\begin{pmatrix}A&B\\C&D\end{pmatrix}\begin{pmatrix}A'&B'\\C'&D'\end{pmatrix}
&=\begin{pmatrix}AA'+BC' & AB'+BD'\\ CA'+DC' & CB'+DD'\end{pmatrix},
\label{eq:leanprimer:fbmul}
\end{align}
named \lean{Matrix.fromBlocks\_transpose} and \lean{Matrix.fromBlocks\_multiply}. Both
are what a student would write by hand; the point of naming them is that \lean{rw} can
apply them mechanically.

\paragraph{Concrete matrices.} The notation \lean{!![0, 1; 1, 0]} builds a literal matrix
over \lean{Fin 2}, and \lean{!![R \^{} 2]} a $1\times1$ matrix over \lean{Fin 1}.

\paragraph{Operations.} \lean{Gᵀ} is the transpose, \lean{Gᴴ} the conjugate transpose
(over $\R$ they coincide, and \lean{TraceBound.lean} needs a lemma,
\lean{Matrix.conjTranspose\_eq\_transpose\_of\_trivial}, to say so, because Mathlib's
positive-definiteness lemmas are stated with \lean{ᴴ}). \lean{G *ᵥ x} is the
matrix--vector product $Gx$, \lean{x ⬝ᵥ y} the dot product $\sum_i x_i y_i$, so
$x^{\mathsf T}Gx$ is \lean{x ⬝ᵥ (G *ᵥ x)}. \lean{G.trace} and \lean{G.det} are the trace
and determinant.

\paragraph{The inverse is total.} \lean{G⁻¹} is defined for \emph{every} square matrix,
by the adjugate formula $G^{-1} = \det(G)^{-1}\operatorname{adj}G$ with the convention
that the ring inverse of a non-unit is $0$. So \lean{G⁻¹ * G = 1} is \emph{not} a theorem
about all $G$; it is \lean{Matrix.nonsing\_inv\_mul G hdet}, which needs a proof
\lean{hdet : IsUnit G.det} that the determinant is invertible in the ground ring. Over a
field this means $\det G\neq0$; over $\Z$ it means $\det G=\pm1$, which is why
unimodularity is exactly the condition for an integer Gram matrix to have an integer
inverse.

\paragraph{Positivity.} Mathlib's definition of positive semidefiniteness, from the
file \lean{PosDef.lean} in \lean{Mathlib/LinearAlgebra/Matrix/}, is
\begin{leancode}
def PosSemidef (M : Matrix n n R) : Prop :=
  M.IsHermitian ∧
    ∀ x : n →₀ R, 0 ≤ x.sum fun i xi ↦ x.sum fun j xj ↦ star xi * M i j * xj
\end{leancode}
that is, $M$ is Hermitian and $\sum_{ij}\bar x_i M_{ij}x_j\ge0$ for all finitely
supported $x$. \lean{PosDef} is defined in the same file by the same two clauses, with the
inequality strict for every $x\neq0$ (its binder for \lean{x} is Lean's ``strict
implicit'' bracket, which we do not reproduce here). A \lean{PosDef} hypothesis therefore contains a symmetry statement
for free, which \lean{dualScale\_ge} extracts as \lean{hG.isHermitian}.

\begin{pitfallbox}[Three ways to be wrong about \lean{G⁻¹}]
(i) Forgetting that \lean{G⁻¹} exists for singular \lean{G} and equals $0$ there: a theorem
about \lean{G⁻¹} with no invertibility hypothesis is usually vacuous or false.
(ii) Expecting \lean{(G⁻¹)⁻¹ = G} to be definitional: it is the lemma
\lean{Matrix.nonsing\_inv\_nonsing\_inv}, and it too needs \lean{IsUnit G.det}.
(iii) Over $\Z$, reading \lean{IsUnit G.det} as $\det G\neq0$: it is $\det G=\pm1$.
\end{pitfallbox}

\begin{example}[An integer inverse as a certificate]\label{ex:leanprimer:certificate}
\index{certificate (proof by witness)}
Let $U=\begin{psmallmatrix}0&1\\1&0\end{psmallmatrix}$, the hyperbolic plane of
\cref{ch:lattices}. Then $UU=\mathbf 1$ and $\det U=-1$. The theorem
\begin{leancode}
theorem isUnimodular_of_mul_eq_one {n : ℕ} (G H : Gram n) (h : H * G = 1) :
    IsUnimodular G := by ...
\end{leancode}
turns the one-line computation $UU=\mathbf 1$ into a proof that $\det U=\pm1$, without
ever computing the determinant: from $HG=\mathbf 1$ and multiplicativity,
$\det G\cdot\det H=\det\mathbf 1=1$ in $\Z$, and the only factorizations of $1$ in $\Z$
are $1\cdot1$ and $(-1)(-1)$. The same theorem is instantiated, never re-proved, for the
$8\times8$ Cartan matrix of $E_8$, where the kernel checks the 64-entry product
$E_8^{-1}E_8=\mathbf 1$ directly. Certificates of this kind---an explicit witness whose
verification is mechanical---are the cheapest kind of Lean proof, and \cref{ch:pipeline}
reports that a small local prover closed essentially all of them
\source{STREAM2\_WORKFLOW.md §6, ll.~107--114}.
\end{example}

%=====================================================================================
\section{The tactic toolbox}\label{sec:leanprimer:tactics}

A proof written as an explicit term is unreadable for anything longer than a line. Instead,
after \lean{:= by}, one writes \emph{tactics}\index{tactic}: commands that transform the
current \emph{goal} (a list of hypotheses and a target proposition) until no goal is left.
Each tactic, behind the scenes, builds a piece of the proof term; the kernel checks the
assembled term at the end. \Cref{tab:leanprimer:tactics} lists every tactic used in the
proofs quoted in this book, with the file in which a student can see it at work. The
remainder of this section shows six of them on real goals.

\begin{table}[ht]\centering\footnotesize
\begin{tabular}{@{}p{2.5cm}p{6.4cm}p{5.7cm}@{}}\toprule
Tactic & What it does to the goal & Seen in \\ \midrule
\lean{unfold f} & replaces \lean{f} by its definition & \lean{thetaShift\_isODD} \\
\lean{rw [h]} & rewrites the goal left-to-right with equation \lean{h}; \lean{← h} goes right-to-left & \lean{isUnimodular\_of\_mul\_eq\_one} \\
\lean{exact t} & closes the goal with the term \lean{t}, whose type must match exactly & \lean{isUnimodular\_of\_mul\_eq\_one} \\
\lean{apply t} & closes the goal with \lean{t} applied to new goals for its missing hypotheses & \lean{circle\_effective\_scale\_ge\_two} \\
\lean{intro x} & moves a $\forall x$ or a hypothesis from the target into the context & \lean{cartanE8\_evenDiag} \\
\lean{rintro ⟨a, h, rfl⟩} & \lean{intro} that also destructures pairs, existentials and equations & \lean{isSection\_image} \\
\lean{constructor} & splits \lean{P ∧ Q} into two goals (or \lean{↔} into two implications) & \lean{index\_mod\_eight} \\
\lean{rcases h with a | b} & case split on an \lean{∨} or an inductive hypothesis & \lean{even\_quadratic\_form\_of\_even\_diag} \\
\lean{have h : P := by ...} & proves an intermediate statement and names it & everywhere \\
\lean{calc} & a chain of equalities or inequalities, each step justified & \lean{isODD\_mul} \\
\lean{ext i j} & to prove two matrices (functions) equal, prove equality at every entry & \lean{cartanE8\_symm} \\
\lean{fin\_cases i} & splits a goal about \lean{i : Fin n} into $n$ concrete goals & \lean{eta\_one\_reindex} \\
\lean{simp [lemmas]} & rewrites with the \lean{@[simp]} lemma set plus the ones listed, until nothing changes & \lean{thetaShift\_isODD} \\
\lean{ring} & proves identities in a commutative ring by normalizing both sides & \lean{tduality\_fixes\_left\_momentum} \\
\lean{field\_simp} & clears denominators, given hypotheses that they are non-zero & \lean{tduality\_fixes\_left\_momentum} \\
\lean{norm\_num} & evaluates numerical expressions and decides numeral (in)equalities & \lean{k3k3\_anomaly} \\
\lean{decide} & evaluates a decidable proposition by computation inside the kernel & \lean{cartanE8\_evenDiag} \\
\lean{rfl} & both sides are the same term after unfolding & \lean{add\_pos}, \lean{eta\_one\_reindex} \\
\lean{linarith} & finds a linear combination of hypotheses that proves a linear (in)equality & \lean{dualScale\_ge} \\
\lean{nlinarith [hints]} & \lean{linarith} after multiplying pairs of hypotheses and hints & \lean{self\_dual\_radius\_unique} \\
\lean{t1 <;> t2} & run \lean{t2} on every goal that \lean{t1} produced & \lean{cartanE8\_symm} \\
\bottomrule\end{tabular}
\caption{The tactics used in the proofs quoted in this book.}\label{tab:leanprimer:tactics}
\end{table}

\subsection{\lean{unfold}, \lean{field\_simp}, \lean{ring}: an algebraic identity}
The proof of \lean{tduality\_fixes\_left\_momentum} is three lines:
\begin{leancode}
theorem tduality_fixes_left_momentum (n w R : ℚ) (hR : R ≠ 0) :
    leftMomentum w n (1 / R) = leftMomentum n w R := by
  unfold leftMomentum
  field_simp
  ring
\end{leancode}
After \lean{unfold} the goal is $w/(1/R) + n\cdot(1/R) = n/R + wR$. \lean{field\_simp}
finds the hypothesis \lean{hR} in the context, multiplies through by $R$ and clears every
fraction, leaving a polynomial identity; \lean{ring} normalizes both sides to the same
polynomial and closes the goal. Without \lean{hR} the tactic would refuse: division by zero
is defined in Lean (as $0$), so $w/(1/R)=wR$ is false at $R=0$, and the kernel would
never accept the theorem without the hypothesis. This is a small example of how the
hypothesis list is forced by the mathematics and not by taste.

\subsection{\lean{constructor}, \lean{rfl}: a conjunction of numerical facts}
\begin{leancode}
theorem index_mod_eight :
    sigK3.index % 8 = 0 ∧ sigMukai.index % 8 = 0 ∧ sigK3T2.index % 8 = 0 := by
  constructor
  · rfl
  constructor
  · rfl
  · rfl
\end{leancode}
\lean{constructor} splits the conjunction; the bullets \lean{·} focus one goal at a time;
\lean{rfl} succeeds because \lean{sigK3.index} unfolds by computation to the numeral
$3-19=-16$, and $-16 \bmod 8$ evaluates to $0$. The theorem does not prove the mod-8
theorem for even unimodular lattices; it checks three instances of it, and its docstring
says so \source{Lattice/K3T2Signature.lean, ll.~176--185}.

\subsection{\lean{ext}, \lean{fin\_cases}, \lean{decide}: finite matrices}
\begin{leancode}
theorem cartanE8_symm : cartanE8ᵀ = cartanE8 := by
  ext i j
  fin_cases i <;> fin_cases j <;> rfl
\end{leancode}
\lean{ext i j} reduces the matrix equation to the entry equation
\lean{cartanE8ᵀ i j = cartanE8 i j} for arbitrary \lean{i j : Fin 8}. \lean{fin\_cases i}
replaces this by eight goals with \lean{i = 0, ..., 7}; the second \lean{fin\_cases}
replaces each by eight more; \lean{rfl} closes all 64 by evaluation. This is the pattern
for every fixed-size statement in the repository (\lean{eta\_one\_reindex},
\lean{tauShift\_dual\_spec}, the $E_8$ inverse check). Where the entries are not literally
identical but merely equal after arithmetic, \lean{decide} or \lean{norm\_num} replaces
\lean{rfl}. The cost grows as $n^2$ goals, each cheap; the benefit is that no idea is
needed.

\subsection{\lean{nlinarith}: an inequality with a hint}
\index{nlinarith@\lean{nlinarith}}\index{self-dual radius!in Lean}
\begin{leancode}
theorem self_dual_radius_unique (R : ℚ) (hR : 0 < R) (hFix : (1 : ℚ) / R = R) :
    R = 1 := by
  have hR' : R ≠ 0 := ne_of_gt hR
  field_simp at hFix
  nlinarith [sq_nonneg (R - 1), sq_nonneg (R + 1)]
\end{leancode}
After \lean{field\_simp at hFix} the hypothesis reads $R^2=1$ (up to the exact form the
tactic chooses). \lean{linarith} treats the monomials $R$, $R^2$, $R^3$ as independent
unknowns and looks for a linear combination of hypotheses proving the goal; it cannot see
that $R^2=1$ and $R>0$ force $R=1$, because that needs a product. \lean{nlinarith} adds
pairwise products of the hypotheses and the supplied hints, then calls \lean{linarith}. A
certificate of the kind it finds: from the hint $(R-1)^2\ge0$ and $R^2=1$,
$0\le R^2-2R+1 = 2-2R$, so $R\le1$; from $R>0$ times the same hint, $0\le R(R-1)^2 =
R^3-2R^2+R$, and $R^3=R\cdot R^2=R$ (the product of $R$ with \lean{hFix}), so
$0\le 2R-2$ and $R\ge1$. Both inequalities together give $R=1$. The second hint
$(R+1)^2\ge0$ is harmless; the physics behind the first is that $R=1$ is the self-dual
radius of \cref{ch:tduality} in units $\ap=1$, and the theorem says it is the
\emph{unique} positive fixed point of $R\mapsto1/R$.

\subsection{\lean{calc}: closure of the duality group}
\begin{leancode}
theorem isODD_mul (g h : Matrix (Charge d) (Charge d) ℤ) (hg : IsODD g)
    (hh : IsODD h) : IsODD (g * h) := by
  unfold IsODD at *
  calc (g * h)ᵀ * eta d * (g * h)
    = hᵀ * (gᵀ * eta d * g) * h := by
        simp [Matrix.transpose_mul, Matrix.mul_assoc]
    _ = hᵀ * eta d * h := by rw [hg]
    _ = eta d := by rw [hh]
\end{leancode}
A \lean{calc} block is the formal version of a chain of equalities on a blackboard. Each
step names its justification. \lean{unfold IsODD at *} unfolds the definition in the goal
\emph{and} in the hypotheses, so that \lean{hg} becomes the equation
\lean{gᵀ * eta d * g = eta d} which \lean{rw} can use. The first step is
$(gh)^{\mathsf T}\eta(gh) = h^{\mathsf T}(g^{\mathsf T}\eta g)h$, which is just
$(gh)^{\mathsf T}=h^{\mathsf T}g^{\mathsf T}$ and associativity; \lean{simp} with those two
lemmas does the regrouping.

%=====================================================================================
\section{One proof, line by line: the \texorpdfstring{$\Theta$}{Theta}-shift}
\label{sec:leanprimer:walkthrough}\index{Theta-shift@$\Theta$-shift!in Lean}
\index{simp@\lean{simp}}

We now read a complete proof at the speed of the kernel. The theorem is the first
generator of $\Odd{d}$ in \cref{ch:odd}: an integer antisymmetric shift of the
$B$-field, $g_\Theta = \begin{psmallmatrix}\mathbf 1&\Theta\\0&\mathbf 1\end{psmallmatrix}$,
preserves the momentum--winding pairing $\eta=\begin{psmallmatrix}0&\mathbf 1\\\mathbf 1&0\end{psmallmatrix}$.
Giveon, Porrati and Rabinovici write the element, require $\Theta_{ij}\in\Z$ and
$\Theta_{ji}=-\Theta_{ij}$, and explain the physics: the constant $B$-term in the action is
a total derivative, so an integer shift changes the action by a multiple of $2\pi$ and drops
out of the path integral \source{GPR §2.4, ll.~1355--1373}.

\subsection{The computation on paper}
With $g=g_\Theta$,
\begin{equation}
g^{\mathsf T}\eta g
=\begin{pmatrix}\mathbf 1&0\\\Theta^{\mathsf T}&\mathbf 1\end{pmatrix}
 \begin{pmatrix}0&\mathbf 1\\\mathbf 1&0\end{pmatrix}
 \begin{pmatrix}\mathbf 1&\Theta\\0&\mathbf 1\end{pmatrix}
=\begin{pmatrix}\mathbf 1&0\\\Theta^{\mathsf T}&\mathbf 1\end{pmatrix}
 \begin{pmatrix}0&\mathbf 1\\\mathbf 1&\Theta\end{pmatrix}
=\begin{pmatrix}0&\mathbf 1\\\mathbf 1&\Theta^{\mathsf T}+\Theta\end{pmatrix}.
\label{eq:leanprimer:thetacomp}
\end{equation}
So $g^{\mathsf T}\eta g=\eta$ exactly when $\Theta^{\mathsf T}+\Theta=0$. (We checked
\eqref{eq:leanprimer:thetacomp} symbolically for $d=2$ with a general $\Theta$; the
residual lower-right block is $\Theta+\Theta^{\mathsf T}$, as it should be. This is a
symbolic check, not part of the Tier A claim.) The paper computation has three steps:
transpose $g$, multiply blocks, use antisymmetry. The Lean proof has the same three steps.

\subsection{The proof}
\begin{leancode}
theorem thetaShift_isODD (Θ : Matrix (Fin d) (Fin d) ℤ) (hΘ : Θᵀ = -Θ) :
    IsODD (thetaShift Θ) := by
  unfold IsODD thetaShift eta
  rw [fromBlocks_transpose, fromBlocks_multiply, fromBlocks_multiply,
    Matrix.transpose_one, Matrix.transpose_zero, hΘ]
  simp
\end{leancode}

\paragraph{The statement.} For every $d$ (implicit, from the file's \lean{variable}
line), every integer $d\times d$ matrix $\Theta$, and a proof \lean{hΘ} that
$\Theta^{\mathsf T}=-\Theta$: the matrix \lean{thetaShift Θ} satisfies \lean{IsODD}. Lean
prints the fully elaborated statement, with every implicit argument and namespace made
visible, on request:
\begin{leancode}
#check @DualScaleStream2.TDuality.thetaShift_isODD
-- ∀ {d : ℕ} (Θ : Matrix (Fin d) (Fin d) ℤ),
--   Θ.transpose = -Θ → IsODD (thetaShift Θ)
\end{leancode}
This is the ``if'' direction only. The converse, that \lean{IsODD (thetaShift Θ)} forces
antisymmetry, is true by \eqref{eq:leanprimer:thetacomp} but is not in the file
(\cref{ex:leanprimer:converse}).

\paragraph{Line 1: \lean{unfold IsODD thetaShift eta}.} The goal
\lean{IsODD (thetaShift Θ)} is replaced by the definitions' bodies:
\begin{leancode}
⊢ (fromBlocks 1 Θ 0 1)ᵀ * fromBlocks 0 1 1 0 * fromBlocks 1 Θ 0 1
    = fromBlocks 0 1 1 0
\end{leancode}
Here \lean{1} and \lean{0} are the $d\times d$ identity and zero blocks; matrix
multiplication associates to the left, so the expression is $(g^{\mathsf T}\eta)g$.

\paragraph{Line 2: the \lean{rw} chain.} Rewrites are applied in order, each to the first
matching sub-expression.
\begin{itemize}[leftmargin=*,itemsep=1pt]
\item \lean{fromBlocks\_transpose}, equation \eqref{eq:leanprimer:fbT}, turns
  \lean{(fromBlocks 1 Θ 0 1)ᵀ} into \lean{fromBlocks 1ᵀ 0ᵀ Θᵀ 1ᵀ}.
\item The first \lean{fromBlocks\_multiply}, equation \eqref{eq:leanprimer:fbmul},
  multiplies the two leftmost block matrices. The goal's left side becomes
  \lean{fromBlocks (1ᵀ * 0 + 0ᵀ * 1) (1ᵀ * 1 + 0ᵀ * 0) (Θᵀ * 0 + 1ᵀ * 1) (Θᵀ * 1 + 1ᵀ * 0)}
  times \lean{fromBlocks 1 Θ 0 1}.
\item The second \lean{fromBlocks\_multiply} performs the remaining product, producing a
  single \lean{fromBlocks} whose four entries are sums of products of \lean{0},
  \lean{1}, \lean{Θ} and \lean{Θᵀ}. The lower-right entry is
  \lean{(Θᵀ * 0 + 1ᵀ * 1) * Θ + (Θᵀ * 1 + 1ᵀ * 0) * 1}.
\item \lean{Matrix.transpose\_one} and \lean{Matrix.transpose\_zero} replace \lean{1ᵀ}
  by \lean{1} and \lean{0ᵀ} by \lean{0}.
\item \lean{hΘ} replaces every \lean{Θᵀ} by \lean{-Θ}. This is the only place the
  hypothesis is used; the lower-right entry now reads
  \lean{(-Θ * 0 + 1 * 1) * Θ + (-Θ * 1 + 1 * 0) * 1}.
\end{itemize}

\paragraph{Line 3: \lean{simp}.} The simplifier applies the ring lemmas for \lean{0} and
\lean{1} (\lean{mul\_zero}, \lean{one\_mul}, \lean{add\_zero}, \lean{neg\_mul}, ...)
until every block is in normal form. The upper-left block becomes \lean{0}, both
off-diagonal blocks \lean{1}, and the lower-right block
\lean{1 * Θ + (-Θ) * 1 = Θ - Θ = 0}. The goal is then
\lean{fromBlocks 0 1 1 0 = fromBlocks 0 1 1 0}, which \lean{simp} closes by reflexivity.
Had \lean{hΘ} been omitted from the \lean{rw} list, \lean{simp} would stop with the
unsolved goal \lean{Θ + Θᵀ = 0}: the same residual as in
\eqref{eq:leanprimer:thetacomp}. That is what ``the proof uses the hypothesis exactly where
the mathematics does'' means in practice.

\paragraph{What the kernel saw.} The tactics produced a proof term whose leaves are the
Mathlib lemmas named above, the hypothesis \lean{hΘ}, and reflexivity. The term type-checks
against the statement; the check depends only on the three standard axioms
(\cref{sec:leanprimer:axioms}). Note that $d$ was never specialized: the proof is one term
valid for all $d$, which is why no \lean{fin\_cases} appears. \Cref{ch:pipeline} records
that this goal was closed by the mid-tier language model, not by the local prover, which
handles only fixed-size goals \source{STREAM2\_WORKFLOW.md §6, ll.~109--114}.

%=====================================================================================
\section{Reading a longer proof: the dual-scale bound}\label{sec:leanprimer:longproof}
\index{dual-scale bound!in Lean}

Most proofs in the repository are longer than three lines, and one reads them by their
\lean{have} skeleton: the intermediate statements are the proof; the tactics between them
are bookkeeping. \lean{dualScale\_ge} in \lean{DualScaleStream2/DualScale/TraceBound.lean}
proves the bound of \cref{ch:dualscale},
\begin{equation}
\Tr G+\Tr G^{-1}\ \ge\ 2d\qquad\text{for every positive-definite }G\in\R^{d\times d},
\label{eq:leanprimer:bound}
\end{equation}
where $\Tr G+\Tr G^{-1}$ is the trace of the generalized metric $\HH(G,0)$ of
\cref{ch:genmetric}, defined in Lean as
\lean{dualScale G := (genMetric G 0).trace} with
\lean{genMetric G B := fromBlocks (G - B * G⁻¹ * B) (B * G⁻¹) (-(G⁻¹ * B)) G⁻¹}.
The idea that avoids diagonalizing $G$ is the identity
\begin{equation}
(G-\mathbf 1)\,G^{-1}\,(G-\mathbf 1) = (G-\mathbf 1)-(\mathbf 1-G^{-1}),
\qquad\text{so}\qquad
\Tr\big[(G-\mathbf 1)G^{-1}(G-\mathbf 1)\big] = \Tr G+\Tr G^{-1}-2d,
\label{eq:leanprimer:sandwich}
\end{equation}
together with the fact that the left side is a trace of $M^{\mathsf T}PM$ with $P=G^{-1}$
positive semidefinite, hence non-negative. (We verified \eqref{eq:leanprimer:sandwich} with
a symbolic $2\times2$ matrix; the expansion is
$(G-\mathbf 1)G^{-1}(G-\mathbf 1)=G-\mathbf 1-\mathbf 1+G^{-1}$ using $GG^{-1}=G^{-1}G=\mathbf 1$.)
The skeleton of the Lean proof is exactly this list:
\begin{leancode}
theorem dualScale_ge (G : Matrix (Fin d) (Fin d) ℝ) (hG : G.PosDef) :
    (2 * d : ℝ) ≤ dualScale G := by
  rw [dualScale_eq]                                  -- goal: 2d ≤ tr G + tr G⁻¹
  have hdet : IsUnit G.det := ...                    -- PosDef gives invertible
  have hl : G⁻¹ * G = 1 := Matrix.nonsing_inv_mul G hdet
  have hr : G * G⁻¹ = 1 := Matrix.mul_nonsing_inv G hdet
  have hGinv_psd : G⁻¹.PosSemidef := hG.posSemidef.inv
  have hGT : Gᵀ = G := ...                           -- from hG.isHermitian
  have hMh : (G - 1)ᴴ = G - 1 := ...
  have hpsd : ((G - 1)ᴴ * G⁻¹ * (G - 1)).PosSemidef :=
    Matrix.PosSemidef.conjTranspose_mul_mul_same hGinv_psd (G - 1)
  rw [hMh] at hpsd
  have htrace_nonneg : 0 ≤ ((G - 1) * G⁻¹ * (G - 1)).trace := hpsd.trace_nonneg
  have hexpand : (G - 1) * G⁻¹ * (G - 1) = (G - 1) - (1 - G⁻¹) := ...
  rw [hexpand] at htrace_nonneg
  rw [Matrix.trace_sub, Matrix.trace_sub, Matrix.trace_sub, Matrix.trace_one]
    at htrace_nonneg
  simp only [Fintype.card_fin] at htrace_nonneg      -- tr 1 = d
  linarith
\end{leancode}
Every \lean{have} is a sentence of the paper proof. The three Mathlib inputs are named:
positive-definite implies invertible (\lean{hdet}); the inverse of a positive-semidefinite
matrix is positive-semidefinite (\lean{PosSemidef.inv}); a sandwich $M^{\mathsf H}PM$ of a
positive-semidefinite $P$ is positive-semidefinite and has non-negative trace (the lemmas
\lean{conjTranspose\_mul\_mul\_same} and \lean{trace\_nonneg} of the \lean{PosSemidef}
namespace). The final \lean{linarith} sees the hypothesis
$0\le(\Tr G-d)-(d-\Tr G^{-1})$ and the goal $2d\le\Tr G+\Tr G^{-1}$, which are the same
linear inequality in the atoms $\Tr G$, $\Tr G^{-1}$, $d$.

\begin{example}[Numbers]\label{ex:leanprimer:numbers}
Take $d=2$ and $G=\diag(4,\tfrac14)$, a rectangular torus with radii $2$ and $\tfrac12$ in
string units. Then $\Tr G+\Tr G^{-1} = 4+\tfrac14+\tfrac14+4 = 8.5\ge4=2d$. The sandwich
is $\diag\big(3^2/4,\ (-\tfrac34)^2\cdot4\big)=\diag(2.25,2.25)$ with trace $4.5$, and
indeed $8.5-4=4.5$. Equality in \eqref{eq:leanprimer:bound} holds only for $G=\mathbf 1$,
where the sandwich vanishes; Lean proves the value $2d$ at $G=\mathbf 1$ as
\lean{dualScale\_one} but not uniqueness of the minimizer.
\end{example}

The circle case is the corollary \lean{circle\_effective\_scale\_ge\_two}:
$2\le R+R^{-1}$ for $R>0$, proved from $R+R^{-1}-2=(R-1)^2/R$ with \lean{field\_simp},
\lean{ring}, \lean{div\_nonneg} and \lean{linarith}. This is the dual-scale inequality of
\cref{ch:dualscale} in units $\ap=1$; its status as mathematics is \TA, its physical
reading is discussed there.

%=====================================================================================
\section{Reading \lean{\#print axioms}}\label{sec:leanprimer:axioms}
\index{print axioms@\lean{\#print axioms}}\index{axioms (Lean)}

A green build is not the whole gate. Two things can hide inside a compiled file. A
\lean{sorry}\index{sorry@\lean{sorry}}\index{native decide@\lean{native\_decide}} is a
placeholder that makes Lean accept any statement; it is reported as a
warning, and a theorem that merely \emph{uses} a sorried lemma inherits the gap silently.
And \lean{native\_decide} asks Lean to trust the compiled evaluator rather than the kernel,
adding the axiom \lean{Lean.ofReduceBool}. The command \lean{\#print axioms} exposes both.
We ran it, from a scratch file that imports the four modules of this chapter, on the
project's built library:
\begin{leancode}
#print axioms DualScaleStream2.TDuality.thetaShift_isODD
-- 'DualScaleStream2.TDuality.thetaShift_isODD' depends on axioms:
--   [propext, Classical.choice, Quot.sound]
#print axioms DualScaleStream2.DualScale.dualScale_ge
-- [propext, Classical.choice, Quot.sound]
#print axioms DualScaleStream2.Lattice.isUnimodular_of_mul_eq_one
-- [propext, Classical.choice, Quot.sound]
#print axioms StringTheory.UseCases.TDuality.self_dual_radius_unique
-- [propext, Classical.choice, Quot.sound]
\end{leancode}
Anything else in the list---\lean{sorryAx}, \lean{Lean.ofReduceBool}, or a user-declared
\lean{axiom}---means the theorem is not Tier A. The project's gate script
\lean{tools/axiom\_audit.py} runs this command over every theorem of a library and fails
unless each depends only on the three standard axioms
\source{STREAM2\_WORKFLOW.md §2, ll.~22--24}. It earned its place: the first run found a
\lean{native\_decide} in a Stream 1 file that every build and every \lean{sorry} search had
passed, and it was replaced by a kernel-checked \lean{norm\_num} proof
\source{LL.md §S2.5, ll.~56--61}.

\begin{pitfallbox}[Zero axioms, one hundred percent]
Every Mathlib-based theorem depends on \lean{propext}, \lean{Classical.choice} and
\lean{Quot.sound}; ``zero axioms'' is not a state a theorem in this book can be in, and a
document that claims it has not run the command. Likewise a count of theorems in a
vendored library is not a measure of what this project proved; the project's own log
records a headline figure of that kind as an overclaim \source{LL.md, ll.~3--13}.
\end{pitfallbox}

%=====================================================================================
\section{Checking a theorem yourself}\label{sec:leanprimer:setup}

The whole point of a machine-checked book is that the reader does not have to believe the
authors. The steps below reproduce the check on a Linux or macOS machine with enough free disk for
Mathlib's compiled object files (several gigabytes); the toolchain and Mathlib version are
pinned in the repository.\index{lake@\lean{lake}}
\begin{enumerate}[leftmargin=*,itemsep=2pt]
\item Install \lean{elan}, the Lean version manager (one shell command from the Lean
  website). It reads the file \lean{lean-toolchain} in the repository, which pins
  \lean{leanprover/lean4:v4.33.1}, and installs exactly that compiler.
\item Clone the repository and run \lean{lake exe cache get} in its root. This downloads
  pre-compiled Mathlib object files for the pinned Mathlib tag \lean{v4.33.1} instead of
  compiling Mathlib, which would take hours; the repository's setup notes record that
  \lean{lake update} triggers this automatically \source{docs/INFRA\_SETUP.md, ll.~53--60}.
\item Run \lean{lake build DualScaleStream2}. The Stream 2 library is a separate
  \lean{lean\_lib} in \lean{lakefile.lean}, built by name so that a failure there can never
  turn the Stream 1 core red \source{STREAM2\_WORKFLOW.md §2, ll.~32--33}. A clean build
  ends with no error and no \lean{sorry} warning.
\item Write a probe file, say \lean{Probe.lean}, containing the imports and the
  \lean{\#print axioms} lines of \cref{sec:leanprimer:axioms}, and run
  \lean{lake env lean -DmaxHeartbeats=1000000 Probe.lean}. The output should be the four
  lines shown there.
\item To read a statement with all implicit arguments, add \lean{\#check @name}. To see a
  goal state, open the file in VS Code with the Lean 4 extension, place the cursor after
  any tactic line, and read the ``Lean Infoview'' panel; this is how the goal states of
  \cref{sec:leanprimer:walkthrough} were obtained.
\end{enumerate}

\begin{pitfallbox}[Heartbeats: the single-file check is stricter than the build]
\index{heartbeats (Lean)}Lean bounds the work per declaration by a ``heartbeat'' budget. The repository's
\lean{lakefile.lean} raises it to $10^6$ for \lean{lake build}, but a bare
\lean{lake env lean file.lean} uses the default $2\times10^5$ and can time out on a proof
that the real build accepts. Always pass \lean{-DmaxHeartbeats=1000000} (and
\lean{-DmaxRecDepth=8000}) to single-file compiles. A stricter budget can reject a valid
proof; it can never accept an invalid one, so no wrong result enters the library this way,
only wasted time \source{LL.md §S2.4, ll.~47--54}.
\end{pitfallbox}

%=====================================================================================
\section{What the machine has checked}\label{sec:leanprimer:lean}

The four files used as examples in this chapter are quoted in full elsewhere in the book
(\cref{ch:lattices,ch:odd,ch:dualscale,ch:tduality}). Here we record only what has been
verified about them from the point of view of this chapter: the statements as read, the
axioms as printed, and the gaps a reader should not paper over. All declarations below
compile in the project's build of 2026-09-17 and depend only on \lean{propext},
\lean{Classical.choice} and \lean{Quot.sound} (\cref{sec:leanprimer:axioms}).

\begin{leanbox}[Lattice predicates and the two general lemmas]
\begin{leancode}
theorem isUnimodular_of_mul_eq_one {n : ℕ} (G H : Gram n) (h : H * G = 1) :
    IsUnimodular G := by ...
theorem even_quadratic_form_of_even_diag {n : ℕ} (G : Gram n) (hsymm : Gᵀ = G)
    (heven : IsEvenDiag G) (x : Fin n → ℤ) : Even (x ⬝ᵥ (G *ᵥ x)) := by ...
\end{leancode}
\textbf{Reading.} For every rank $n$ and integer Gram matrix $G$: an integer left inverse
proves $\det G=\pm1$; and if $G$ is symmetric with even diagonal then $x^{\mathsf T}Gx$ is
even for every integer vector $x$. \textbf{Proof idea.} The first is
$\det G\det H=1$ in $\Z$ (\lean{Matrix.det\_mul}, \lean{Matrix.det\_one},
\lean{Int.eq\_one\_or\_neg\_one\_of\_mul\_eq\_one}). The second splits
$G=D+U+U^{\mathsf T}$ into diagonal and strict upper triangle by a three-way case split
\lean{rcases lt\_trichotomy i j}, then $x^{\mathsf T}Gx = x^{\mathsf T}Dx + 2\,x^{\mathsf T}Ux$.
\textbf{Not proved.} The converse (an even form has even diagonal) is not stated; no
abstract lattice, isomorphism or classification appears
\source{Lattice/Basic.lean header, ll.~40--45}. \TA
\end{leanbox}

\begin{leanbox}[The duality group for every $d$]
\begin{leancode}
theorem thetaShift_isODD (Θ : Matrix (Fin d) (Fin d) ℤ) (hΘ : Θᵀ = -Θ) :
    IsODD (thetaShift Θ) := by ...
theorem basisChange_isODD (A B : Matrix (Fin d) (Fin d) ℤ) (h : Aᵀ * B = 1) :
    IsODD (basisChange A B) := by ...
theorem eta_mul_self : eta d * eta d = 1 := by ...
theorem eta_isODD : IsODD (eta d) := by ...
theorem isODD_mul (g h : Matrix (Charge d) (Charge d) ℤ) (hg : IsODD g)
    (hh : IsODD h) : IsODD (g * h) := by ...
theorem eta_one_reindex :
    (eta 1).reindex finSumFinEquiv finSumFinEquiv
      = DualScaleStream2.Lattice.hyperbolicU := by ...
\end{leancode}
\textbf{Reading.} The $\Theta$-shift with antisymmetric $\Theta$, the basis change
$\diag(A,B)$ with $A^{\mathsf T}B=\mathbf 1$, and $\eta$ itself satisfy
$g^{\mathsf T}\eta g=\eta$; the predicate is closed under products; at $d=1$ the matrix
$\eta$, reindexed from \lean{Fin 1 ⊕ Fin 1} to \lean{Fin 2}, is literally the hyperbolic
plane $U$. Each of the first two is one direction of an ``iff''. \textbf{Proof idea.}
Block transpose, two block multiplications, the hypothesis, \lean{simp}
(\cref{sec:leanprimer:walkthrough}). \textbf{Not proved.} That the three families generate
$\Odd{d}$; any group structure; the moduli-space quotient
\source{TDuality/ODD.lean header, ll.~33--36}. \TA
\end{leanbox}

\begin{leanbox}[The dual-scale bound]
\begin{leancode}
noncomputable def dualScale (G : Matrix (Fin d) (Fin d) ℝ) : ℝ :=
  (genMetric G 0).trace
theorem dualScale_eq (G : Matrix (Fin d) (Fin d) ℝ) :
    dualScale G = G.trace + G⁻¹.trace := by ...
theorem dualScale_inv (G : Matrix (Fin d) (Fin d) ℝ) (hG : IsUnit G.det) :
    dualScale G⁻¹ = dualScale G := by ...
theorem dualScale_ge (G : Matrix (Fin d) (Fin d) ℝ) (hG : G.PosDef) :
    (2 * d : ℝ) ≤ dualScale G := by ...
theorem dualScale_one :
    dualScale (1 : Matrix (Fin d) (Fin d) ℝ) = 2 * d := by ...
theorem dualScale_circle (R : ℝ) (hR : R ≠ 0) :
    dualScale (!![R ^ 2] : Matrix (Fin 1) (Fin 1) ℝ)
      = (R + R⁻¹) ^ 2 - 2 := by ...
theorem circle_effective_scale_ge_two (R : ℝ) (hR : 0 < R) :
    2 ≤ R + R⁻¹ := by ...
\end{leancode}
\textbf{Reading.} Over the real numbers, for every $d$: the trace of $\HH(G,0)$ is
$\Tr G+\Tr G^{-1}$; it is invariant under $G\mapsto G^{-1}$ for invertible $G$; it is at
least $2d$ for positive-definite $G$; it equals $2d$ at $G=\mathbf 1$; at $d=1$ with
$G=R^2$ it equals $(R+1/R)^2-2$; and $R+1/R\ge2$ for $R>0$. \lean{noncomputable} marks a
definition involving real numbers, which cannot be evaluated by a program; it does not
weaken the theorems. \textbf{Proof idea.} \Cref{sec:leanprimer:longproof}. \textbf{Not
proved.} That $G=\mathbf 1$ is the unique minimizer; anything at $B\neq0$; any physical
interpretation, which the file itself labels Tier C
\source{DualScale/TraceBound.lean header, ll.~44--47}. \TA
\end{leanbox}

\begin{leanbox}[The circle mass spectrum over $\Q$]
\begin{leancode}
theorem tduality_fixes_left_momentum (n w R : ℚ) (hR : R ≠ 0) :
    leftMomentum w n (1 / R) = leftMomentum n w R := by ...
theorem tduality_flips_right_momentum (n w R : ℚ) (hR : R ≠ 0) :
    rightMomentum w n (1 / R) = - rightMomentum n w R := by ...
theorem tduality_invariant_mass_squared (n w R : ℚ) (hR : R ≠ 0) :
    momentumMassSq w n (1 / R) = momentumMassSq n w R := by ...
theorem self_dual_radius_unique (R : ℚ) (hR : 0 < R) (hFix : (1 : ℚ) / R = R) :
    R = 1 := by ...
\end{leancode}
\textbf{Reading.} With $\ap=1$ and $n,w,R$ rational: T-duality fixes $p_L$, negates
$p_R$, hence fixes $p_L^2+p_R^2$; and $R=1$ is the only positive fixed point of
$R\mapsto1/R$. \textbf{Proof idea.} \lean{unfold}, \lean{field\_simp}, \lean{ring};
\lean{nlinarith} with two square hints (\cref{sec:leanprimer:tactics}). \textbf{Not
proved.} Nothing about oscillators or the full CFT; $n$ and $w$ are not restricted to
integers, so the theorem is about rational functions, an arithmetic shadow of the physical
statement whose algebra happens to be identical \source{TDualityMassSpectrum.lean header,
ll.~35--50}. \TA
\end{leanbox}

\begin{center}\small
\begin{tabular}{@{}p{5.6cm}cp{7.2cm}@{}}\toprule
Claim & Tier & Where \\ \midrule
The theorems above, as stated & \TA & the four files; axioms printed in \cref{sec:leanprimer:axioms} \\
$g_\Theta\in\Odd{d}$ \emph{iff} $\Theta$ antisymmetric & \TL & GPR §2.4; ``only if'' by \eqref{eq:leanprimer:thetacomp}, not in Lean \\
$g_\Theta$ is a symmetry of the string spectrum & \TL & GPR §2.4 (total-derivative argument) \\
\lean{leftMomentum} \emph{is} the physical $p_L$ & --- & an identification; never Tier A \\
\bottomrule\end{tabular}
\end{center}

%=====================================================================================
\begin{summarybox}
\begin{itemize}[leftmargin=*,itemsep=1pt]
\item A theorem is a type in \lean{Prop}; a proof is a term of that type; the kernel checks
  the term against the type and trusts nothing else. Three things remain to be trusted: the
  kernel, the three axioms \lean{propext}, \lean{Classical.choice}, \lean{Quot.sound}, and
  the statement itself.
\item \lean{def} makes an object that can be unfolded; \lean{theorem} makes a proof that can
  only be applied. Read the hypothesis list first: it tells you what a formal theorem
  covers.
\item A Mathlib matrix is a function of two indices; \lean{fromBlocks} builds
  $2\times2$ block matrices on \lean{Fin d ⊕ Fin d}; \lean{G⁻¹} is total and needs
  \lean{IsUnit G.det} to behave; \lean{PosDef} includes symmetry.
\item Tactics transform goals. Identities: \lean{unfold}, \lean{rw}, \lean{simp},
  \lean{ring}, \lean{field\_simp}. Finite computations: \lean{ext}, \lean{fin\_cases},
  \lean{decide}, \lean{norm\_num}, \lean{rfl}. Inequalities: \lean{linarith},
  \lean{nlinarith}. Structure: \lean{have}, \lean{calc}, \lean{constructor},
  \lean{intro}.
\item \lean{thetaShift\_isODD} is proved by unfolding, two block multiplications, one use
  of antisymmetry and \lean{simp}; the hypothesis is used exactly where the paper
  computation uses it. It proves one direction only.
\item \lean{\#print axioms} is the last gate: \lean{sorryAx} or \lean{Lean.ofReduceBool}
  in the list disqualifies a theorem from Tier A. Single-file checks need
  \lean{-DmaxHeartbeats=1000000}.
\end{itemize}
\end{summarybox}

%=====================================================================================
\section*{Exercises}
\addcontentsline{toc}{section}{Exercises}

\begin{exercise}[$\star$ Reading statements]\label{ex:leanprimer:reading}
For each of \lean{dualScale\_inv}, \lean{dualScale\_ge} and \lean{eta\_one\_reindex},
write in one English sentence what is quantified, what is assumed, and what is concluded.
Then say which of the three would still be true, and which would become false or
meaningless, if its hypothesis were deleted.
\end{exercise}

\begin{exercise}[$\star$ The other generator]\label{ex:leanprimer:basischange}
Repeat the paper computation \eqref{eq:leanprimer:thetacomp} for
$g=\diag(A,B)$ and show that $g^{\mathsf T}\eta g=\eta$ is equivalent to
$A^{\mathsf T}B=\mathbf 1$ and $B^{\mathsf T}A=\mathbf 1$. Explain why the Lean proof of
\lean{basisChange\_isODD} needs the auxiliary \lean{have h' : Bᵀ * A = 1}, and how it
obtains it from \lean{h} by transposing.
\end{exercise}

\begin{exercise}[$\star$ Goal states]\label{ex:leanprimer:goals}
Open \lean{ODD.lean} in an editor with the Lean extension, place the cursor after the
\lean{rw} line of \lean{thetaShift\_isODD}, and copy the goal. Identify the four blocks and
match them with the entries of \eqref{eq:leanprimer:thetacomp}. Then delete \lean{hΘ}
from the \lean{rw} list and report the goal \lean{simp} leaves open.
\end{exercise}

\begin{exercise}[$\star\star$ Lean: the converse]\label{ex:leanprimer:converse}
State and prove
\begin{leancode}
theorem thetaShift_isODD_iff (Θ : Matrix (Fin d) (Fin d) ℤ) :
    IsODD (thetaShift Θ) ↔ Θᵀ = -Θ := by ...
\end{leancode}
Use \lean{constructor}; the reverse direction is
\lean{thetaShift\_isODD}. For the forward direction, unfold and rewrite as in
\cref{sec:leanprimer:walkthrough} to reach an equation between two \lean{fromBlocks}
matrices, then extract the lower-right block (Mathlib has \lean{toBlocks₂₂} and the
lemma \lean{toBlocks\_fromBlocks₂₂} in the \lean{Matrix} namespace) and solve for
\lean{Θᵀ}.
\end{exercise}

\begin{exercise}[$\star\star$ Lean: a certificate]\label{ex:leanprimer:certificate2}
Define \lean{U : Gram 2 := !![0, 1; 1, 0]} and prove \lean{IsUnimodular U} by applying
\lean{isUnimodular\_of\_mul\_eq\_one} with \lean{H := U}; prove the side condition
\lean{U * U = 1} by \lean{ext i j; fin\_cases i <;> fin\_cases j <;> rfl} (you may need
\lean{simp [Matrix.mul\_apply, Fin.sum\_univ\_succ]} before \lean{rfl}). Then prove
\lean{IsEvenDiag U}. Compare with \lean{hyperbolicU} in \lean{Lattice/Hyperbolic.lean}.
\end{exercise}

\begin{exercise}[$\star\star$ Lean: numbers]\label{ex:leanprimer:numbers2}
Prove \lean{example : (2 : ℝ) ≤ 4 + (4 : ℝ)⁻¹} with \lean{norm\_num}, and
\lean{example : (2 : ℝ) ≤ dualScale (!![(2 : ℝ) \^{} 2] : Matrix (Fin 1) (Fin 1) ℝ)} by
rewriting with \lean{dualScale\_circle} and then \lean{norm\_num}. Which hypothesis of
\lean{dualScale\_circle} must you discharge, and how?
\end{exercise}

\begin{exercise}[$\star\star$ The sandwich]\label{ex:leanprimer:sandwich}
Verify \eqref{eq:leanprimer:sandwich} by hand for a general invertible $G$, and
compute both sides for $G=\begin{psmallmatrix}2&1\\1&1\end{psmallmatrix}$. Then show that
$\Tr[(G-\mathbf 1)G^{-1}(G-\mathbf 1)]=0$ forces $G=\mathbf 1$ when $G$ is positive
definite, which is the uniqueness statement that \lean{TraceBound.lean} does not contain.
\end{exercise}

\begin{exercise}[$\star\star\star$ Lean: a faithful version]
\label{ex:leanprimer:faithful}
The circle theorems of \cref{sec:leanprimer:types} live over $\Q$ with $\ap=1$. Restate
the invariance theorem \lean{tduality\_invariant\_mass\_squared} over $\R$, with $\ap$ a
positive real parameter: define $p_L = n/R + wR/\ap$ and $p_R = n/R - wR/\ap$ with
$n, w : \Z$ cast to $\R$, and prove that $(n,w,R)\mapsto(w,n,\ap/R)$ fixes $p_L^2+p_R^2$. You will need
\lean{hα : 0 < α'} and \lean{hR : 0 < R} for \lean{field\_simp}. Then add the oscillator
term $(2/\ap)(N+\tilde N-2)$ as a further argument and explain why the proof does not
change.
\end{exercise}

\begin{exercise}[$\star\star\star$ Audit]\label{ex:leanprimer:audit}
Write a file that proves \lean{example : (10 : ℕ) \^{} 3 = 1000} twice, once with
\lean{decide} and once with \lean{native\_decide}, name the two proofs, and run
\lean{\#print axioms} on both. Explain from the output why the project's gate rejects the
second and what would have to be trusted to accept it.
\end{exercise}

%=====================================================================================
\section*{Further reading}
\addcontentsline{toc}{section}{Further reading}
\begin{itemize}[leftmargin=*,itemsep=2pt]
\item The Lean sources used in this chapter: \lean{Basic.lean}, \lean{ODD.lean} and
 \lean{TraceBound.lean} in the directories \lean{Lattice/}, \lean{TDuality/} and
 \lean{DualScale/} of \lean{DualScaleStream2/}, and \lean{TDualityMassSpectrum.lean} in
 \lean{StringTheoryFormalization/UseCases/}. Each header
 states its physical background, its mathematical content, what is not proved, and the
 proof techniques; read the header before the code.
\item The physics of the $\Theta$-shift and the basis change, with the total-derivative
 argument: \source{GPR §2.4, ll.~1355--1373}.
\item The project's gate: the kernel as the only accept criterion, the three standard
 axioms, and ``producer $\neq$ verifier'' \source{STREAM2\_WORKFLOW.md §2, ll.~20--38};
 measured results of the tiered proving pipeline
 \source{STREAM2\_WORKFLOW.md §6, ll.~105--124}. The lessons on statement locking, axiom
 auditing, heartbeats and the structure-field parser pitfall
 \source{LL.md §S2.4--S2.13, ll.~47--112} and \source{LL.md §0.4, ll.~188--199}. The
 architecture of the libraries is described in \cref{ch:architecture}; the atlas of
 theorems in \cref{ch:atlas}; the pipeline in \cref{ch:pipeline}.
\item For Lean itself, the official documentation and Mathlib's reference pages are the
 authoritative sources; the definitions of \lean{Matrix}, \lean{fromBlocks},
 \lean{PosDef} and the nonsingular inverse quoted above are from the pinned Mathlib tag
 \lean{v4.33.1} in \lean{lakefile.lean}. Standard; not checked against a source in this
 book's library.
\end{itemize}
